\section{RISC-V}

O RISC-V é uma \gls{isa} baseada nos princípios da arquitetura \acrshort{risc} (\acrlong{risc}). Diferentemente de arquiteturas proprietárias,
o RISC-V possui uma especificação aberta e livre de royalties, o que permite sua utilização tanto em contextos acadêmicos quanto em aplicações
comerciais.

Essa \gls{isa} define um conjunto básico de instruções, responsável pelas funcionalidades essenciais do processador. No entanto, sua
principal característica é a modularidade, permitindo que diferentes extensões sejam incorporadas conforme as necessidades da aplicação.
Dessa forma, é possível adaptar a arquitetura para cenários variados, desde sistemas embarcados com restrições de recursos até sistemas de
alto desempenho.

\subsection{Extensões}

\begin{description}

\item[I]
Instruções básicas para operações com números inteiros. Conforme descrito em \cite{RISCVUnprivilegedISA}, essa extensão deve estar presente em
qualquer implementação RISC-V, pois define operações fundamentais de aritmética e manipulação de registradores.

\item[Zicsr]
Adiciona suporte às instruções de acesso e manipulação de \glspl{csr}.

\item[H]
Adiciona suporte em hardware à virtualização. Porém, apesar de existir, é muito recente e sua implementação em hardware comercial ainda é limitada.

\item[M]
Adiciona suporte à operações de multiplicação e divisão de números inteiros.

\item[A]
Adiciona suporte à operações atômicas em memória.

\item[F]
Adiciona suporte a números em ponto flutuante de precisão simples (32 bits).

\item[D]
Adiciona suporte a números em ponto flutuante de precisão dupla (64 bits).

\item[C]
Adiciona suporte à instruções comprimidas de 16 bits.

\item[S]
Adiciona suporte ao modo supervisor.

\item[G]
Aglomerado geral de extensões que inclui I, M, A, F, D e C.

\end{description}

Para o escopo deste trabalho, serão consideradas apenas as extensões Zicsr, S e G. As estratégias adotadas para mitigar a ausência da extensão H
serão discutidas ao longo do texto.

\subsection{Níveis de Privilégio}
\subsection{Control and Status Registers (CSRs)}
